\documentclass[12pt]{article}
\usepackage{it_hw}
\input{macro}

\excludeversion{problem}
\includeversion{solution}

\begin{document}

\begin{center}
{\bf Information Theory and Its Application in Deep Learning: Midterm \begin{solution}Solutions\end{solution}} \\ \medskip
\begin{problem}Thursday, April 24, 2025 \end{problem}
\end{center}

\begin{problem}
\noindent 
\begin{itemize}
\item You can use previous parts of a problem even if you did not solve them.
\item As throughout the course, entropy ($H$) and Mutual Information ($I$) are in bits.
\item $\log$ is taken in base 2. 
\item `prefix code' refers to a variable length code satisfying the prefix condition.  
\end{itemize}
\LeftFoot{Midterm}
\end{problem}

\begin{solution}
\LeftFoot{Midterm Solution}
\end{solution}

\begin{problem}

    \vskip .3in
    {\bf ID : \underline{\hspace{6cm}}}\\
    \vskip .3in
    {\bf Name : \underline{\hspace{5.5cm}}}\\
    ~
    \vskip 2in
    \begin{table}[h]
        \Large
        \begin{tabular}{c|c}
           1 &\textcolor{white}{aaaaaaaa}/~30\\
           2 &\textcolor{white}{aaaaaaaa}/~30 \\
           3 &\textcolor{white}{aaaaaaaa}/~40 \\
            \hline
           Total &  \textcolor{white}{aaaaaaaa}/~100 \\
        \end{tabular}
    \end{table}
    \newpage
    {\bf extra sheet}
    \newpage
\end{problem}

\begin{enumerate}
%%% Problem 1: Vin's Idea
%%%%%%%%%%%%%%%%%%%%%
\item  {\bf Albert's Idea} {\it (30 points)}\\
Albert is very excited about a new lossless compression, and claims it can beat entropy.\\

Albert: ``Suppose $X_1,X_2,\ldots$ is an  i.i.d. Bernoulli-$1/2$ sequence.
We can break up this sequence into its pattern of `repeats'.  For instance,
0001100001... begins with repeats (also known as `run-lengths')  `000', `11', and `0000'. If we let $L_i$ be
the length of the $i$th repeat, we can represent the sequence 
by  $(X_1, L_1, L_2, \ldots )$. For example, 
\begin{itemize}
	\item $1010\ldots$ would be represented by $(1,1,1,1, \ldots)$
	\item $11100111110\ldots$ by $(1,3,2,5, \ldots)$ and 
	\item $001011111110\ldots$ by $(0,2,1,1,7, \ldots)$. 
\end{itemize}
In particular, I suggest we  describe the sequence $X_1, X_2, \ldots, X_{\sum_{i=1}^{10} L_i}$ by describing\\
 $(X_1, L_1, L_2,\ldots,L_{10})$, which I'm sure would be a heavily compressed representation!!''

\begin{enumerate}
\item {\it (5 points)} What is the entropy of the first repeat length $H(L_1)$?

\item {\it (5 points)} Describe an optimal prefix code for $L_1$. What is its expected code-length?  

%\item {\it (6 points)} What is the joint entropy of all the repeat lengths $H(L_1,L_2,\ldots,L_{10})$?

\item {\it (5 points)} What is $H(X_1,L_1,\ldots,L_{10})$?

\item {\it (5 points)} Describe an optimal  uniquely decodable code for $(X_1,L_1,\ldots,L_{10})$.   What is its expected code-length?  
Call it ``Albert's code''. 

\item {\it (5 points)} What is the expected number of source symbols $E[\sum_{i=1}^{10} L_i]$ that Albert's code encodes?

\item {\it (5 points)} Comment, based on your answers to the previous two parts, on whether Albert's code is  ``beating entropy'' on average.

%\item {\it (6 points)} Now suppose that $\{X_1, X_2, \ldots\}$ are distributed as i.i.d. Bernoulli-$p$ random
%variables. What is the expected number of bits that Vinith encodes $E[\sum_{i=1}^{10} L_i]$?

\end{enumerate}


%%% Solution of  Problem 1: Vin's Idea
%%%%%%%%%%%%%%%%%%%%%
\vskip 0.2in
\begin{solution}
{\bf Solution:}\\
\begin{enumerate}
\item At each time $i$, the current repeat will end with probability $1/2$. 
Therefore, each repeat length $L_i$ is a geometric random variable with parameter $1/2$. Therefore, 
\[ H(L_1) = \sum_{j=1}^\infty 2^{-j} \log(2^j) = \sum_{j=1}^\infty j2^{-j} = 2 \mbox{.} \]

\item An optimum prefix code for $L_1$ is

\begin{tabular}{cc}
$L_1$ & Codeword  \\
1 & 1 \\
2 & 01\\
3 & 001\\
\vdots & \vdots
\end{tabular}
where its expected code-length is
\[\sum_{i=1}^{\infty} \frac{1}{2^i} i = 2.\]

\item Since the process is memoryless, the repeat lengths $L_i$ are independent and identically distributed. Therefore $H(L_1,\ldots,L_{10}) = 10 H(L_1) = 20$. 
The first symbol $X_1$ is independent of the repeat lengths $(L_1,\ldots,L_{10})$, so the joint entropy is given by
\[
H(X_1, L_1,\ldots,L_{10}) = H(X_1) + 10H(L_1) = 21 \mbox{.}
\]

\item We have to use one bit to describe $X_1$. Since $L_1,L_2,\ldots,L_{10}$ are i.i.d. geometric random variables with parameter $1/2$, we can use the prefix code for $L_1$ (which we did in (b)) repeatedly. Since the expected code-length of an optimal prefix code for $L_1$ is $2$, the expected code-length will be 
\[
1+2\times10 = 21.
\]

\item The expected value of $L_i$ is that of a geometric random variable with parameter
$1/2$: $E[L_i] = 2$. Using the linearity of expectation,
\[
E[\sum_{i=1}^{10} L_i] = \sum_{i=1}^{10} E[L_i] = 20 \mbox{.}
\]

\item Albert is encoding 20 souce bits using an average of 21-bit-long binary codewords. Thus, on average he is expending more than one bit of description per source bit, and not `beating the entropy'. .

\end{enumerate}
\end{solution}

\begin{problem}
    \clearpage
    ~
    \clearpage
\end{problem}

%%% Problem 2: Non-i.i.d. Source
%%%%%%%%%%%%%%%%%%%%%
\vskip 0.3in
\item {\bf Non-i.i.d. Source} {\it (30 points)}\\
Consider a second-order binary Markov process   $\{X_i\}_{i \geq 1}$ characterized as follows:
\begin{itemize}
\item $P(X_1=0,X_2=0) = P(X_1=1,X_2=1) = \frac{1}{6}$ and  \\ $P(X_1=0,X_2=1) = P(X_1=1,X_2=0) = \frac{1}{3}$.
\item For $n\geq 3$,
\begin{itemize}
\item If $X_{n-1}=X_{n-2}$, then $X_n = 1- X_{n-1}$.
\item If $X_{n-1}\neq X_{n-2}$, then $X_n$ is drawn as a fair coin flip,  independent of $\{X_i\}_{i=1}^{n-1}$.
\end{itemize}
\end{itemize}

\begin{enumerate}[(a)]

\item {\it (6 points)} Find an optimal prefix code for the pair $(X_1, X_2)$, along with its expected code-length. 

\item {\it (6 points)} Show that the distribution of $(X_{n}, X_{n+1})$ is the same for all $n \geq 1$ (and, hence, the process is stationary). 


%\item {\it (5 points)} TAKE FROM HERE. Show that $Pr(X_n=0) = Pr(X_n=1) = 1/2$ for all $n\geq 1$.
%\item {\it (5 points)} Let $Y_n = (X_{2n-1}, X_{2n})$, for example $Y_1 = (X_1,X_2)$, $Y_2 = (X_3,X_4)$, and so on. Find an optimal prefix code for 


\item {\it (6 points)} Find the ``entropy rate'' of the process 
\bea
\lim_{n\rightarrow\infty} \frac{H(X_1,X_2,\ldots,X_n)}{n} .
\eea
[Hint: can justify and use the facts that  $H(X_1,X_2,\ldots,X_n) = \sum_{i=1}^n H(X_i | X^{i-1})$, and $H(X_i | X^{i-1}) = H(X_i | X_{i-1}, X_{i-2}) = H(X_3|X_2, X_1)$ for $i \geq 3$ ]


\item {\it (6 points)} For FIXED $n \geq 2$, does there exist a uniquely decodable code for $(X_1,X_2,\ldots,X_n)$ 
whose expected code-length is exactly $H(X_1,X_2,\ldots,X_n)$? If so, describe one. If not, explain why. 

\item {\it (6 points)} Describe a uniquely decodable code (other than arithematic encoding) for $(X_1,X_2,\ldots,X_n)$ 
that attains the entropy rate. That is, a code with length function $\ell_n$ such that 
\bea
\lim_{n\rightarrow\infty} \frac{E \left[  \ell_n (X_1,X_2,\ldots,X_n) \right]}{n} 
\eea
is equal to the entropy rate from part (c). 

\end{enumerate}


%%% Solution of  Problem 2: Non-i.i.d. Source
%%%%%%%%%%%%%%%%%%%%%
\vskip 0.2 in
\begin{solution}
{\bf Solution: }\\
\begin{enumerate}[(a)]
\item The Huffman tree for $(X_1,X_2)$ is

\begin{tabular}{llccccccc}
Codeword & $(X_1,X_2)$\\
0 & $(0,1)$ & $\frac{1}{3}$ & $\frac{1}{3}$ & $\frac{1}{3}$ & 1 \\
10 & $(1,0)$ & $\frac{1}{3}$ & $\frac{1}{3}$ & $\frac{2}{3}$ \\
110 & $(0,0)$ & $\frac{1}{6}$ & $\frac{1}{3}$ \\
111 & $(1,1)$ & $\frac{1}{6}$  \\
\end{tabular}

Its expected code-length is 
\[
1\times\frac{1}{3} + 2\times\frac{1}{3}+3\times\frac{1}{6}+3\times\frac{1}{6} = 2.
\]

\item We will show that $(X_{n},X_{n+1})$ has the same distribution with $(X_1,X_2)$ using induction. Clearly, the statement is true for $n=1$. Suppose $(X_{k-1},X_{k})$ has the same distribution with $(X_1,X_2)$. Then,
\begin{align*}
P(X_k=0,X_{k+1}=0) =& P(X_{k-1}=0, X_k=0,X_{k+1}=0)+P(X_{k-1}=1, X_k=0,X_{k+1}=0) \\
=& P(X_{k+1}=0|X_{k-1}=1, X_k=0)P(X_{k-1}=1, X_k=0) \\
=& \frac{1}{2}\cdot \frac{1}{3}\\
=& \frac{1}{6}\\
P(X_k=0,X_{k+1}=1) =& P(X_{k-1}=0, X_k=0,X_{k+1}=1)+P(X_{k-1}=1, X_k=0,X_{k+1}=1) \\
=& P(X_{k+1}=1|X_{k-1}=0, X_k=0)P(X_{k-1}=0, X_k=0)\\
& +P(X_{k+1}=1|X_{k-1}=1, X_k=0)P(X_{k-1}=1, X_k=0)\\
=& 1\cdot \frac{1}{6} + \frac{1}{2}\cdot \frac{1}{3}\\
=& \frac{1}{3}
\end{align*}
Similarly, it is easy to show that $P(X_k=1,X_{k+1}=0)  = \frac{1}{3}$ and $P(X_k=1,X_{k+1}=1) =\frac{1}{6}$.

\item For $n\geq 3$, $H(X_1,X_2,\ldots,X_n) = (n-2)H(X_3|X_2,X_1) + H(X_1,X_2)$. Therefore,
\bea
\lim_{n\rightarrow\infty} \frac{H(X_1,X_2,\ldots,X_n)}{n} = H(X_3|X_2,X_1).
\eea
Note that the conditional entropy  $H(X_3|X_2,X_1)$ is
\begin{align*}
H(X_3|X_2,X_1) =& \frac{1}{6} \cdot H(X_3|X_2=0,X_1=0) + \frac{1}{6} \cdot H(X_3|X_2=1,X_1=1) \\
&+ \frac{1}{3} \cdot H(X_3|X_2=1,X_1=0) + \frac{1}{3} \cdot H(X_3|X_2=0,X_1=1) \\
=&  \frac{1}{6} \cdot 0 + \frac{1}{6} \cdot 0 + \frac{1}{3} \cdot 1 + \frac{1}{3} \cdot 1\\
=& \frac{2}{3}.
\end{align*}

\item For $n\geq 2$,
\[
P(x_1,x_2,\ldots,x_n) = P(x_1,x_2) \prod_{i=3}^n P(x_i|x_{i-1},x_{i-2})
\]
where $P(x_1,x_2)$ is either $\frac{1}{3}$ or $\frac{1}{6}$, and $P(x_i|x_{i-1},x_{i-2})$ takes value from $\{1, 0, \frac{1}{2}\}$. Therefore, it is not a diadic distribution, we can not achieve $H(X_1,X_2,\ldots,X_n)$.

\item Consider the following coding scheme.
\begin{itemize}
\item Use any code for $X_1,X_2$ (e.g. Huffman). This will be negligible in terms of average code-length.
\item For $n\geq 3$, if $X_{n-1}\neq X_{n-2}$ describe $X_n$ using 1 bit. If $X_{n-1}=X_{n-2}$, then send nothing since $X_n$ will be deterministic.
\end{itemize}
Since $P(X_{n-1}=X_{n-2}) = \frac{2}{3}$, the average code-length per symbol will be $\frac{2}{3}$.

Note: we can use Albert's code to achieve the entropy rate. Let $L_i$ be a length of $i$-th repeat. Then, $Z_i = L_i-1$ is i.i.d. Bernoulli-1/2 random process and $(X_1, Z_1,Z_2,\ldots)$ will be a compressed version of $(X_1,X_2,\ldots)$. For given compressed version $(X_1,Z_1,Z_2,\ldots,Z_m)$, the expected number of encoded source symbols is 
\bea
E[L_1+L_2+\cdots+L_m] = \frac{3}{2}m.
\eea
Therefore, we can argue that 
\[\lim_{n\rightarrow\infty}\frac{E[l_n(X_1,X_2,\ldots,X_n)]}{n} =  \frac{2}{3}.\]
\end{enumerate}
\end{solution}

\begin{problem}
    \clearpage
    ~
    \clearpage
\end{problem}

%%% Problem 3: Entropy of Sum and Difference of I.I.D. Random Variables
%%%%%%%%%%%%%%%%%%%%%
\vskip 0.3in
\item {\bf Entropy of a Sum and a Difference of I.I.D. Random Variables} {\it (40 points)}\\
We will prove that if $(Y,Y')$ are i.i.d. discrete random variables then:
\bea
H(Y-Y')-H(Y)\le 2(H(Y'+Y)-H(Y)).
\eea
We will prove this inequality in the following steps.
\begin{enumerate}
\item \underline{\it Data Processing Inequality for Mutual Information}  {\it (10 points)}\\
Let $X_1,X_2$ be discrete random variables. Also, let $Y_1=F(X_1)$ and $Y_2=G(X_2)$ for some functions, $F(\cdot), G(\cdot)$. Prove that:
\bea
I(X_1;X_2)\ge I(Y_1;Y_2).
\eea

\item \underline{\it Submodularity}  {\it (10 points)}\\

Suppose that there exist functions $F$, $G$ and $R$ such that  $X_0=F(X_1)=G(X_2)$ and $X_{12}=R(X_1,X_2)$, where  $X_1,X_2$ are discrete random variables. 
Use the previous part to prove that
\bea
H(X_{12})+H(X_0)\le H(X_1)+H(X_2).
\eea
%[Hint: Use part (a), with $Y_1 = Y_2 = X_0$.]

\item \underline{\it Ruzsa Triangle Inequality}  {\it (10 points)}\\
Let $X,Y,Z$ be independent discrete random variables.  Use Part (b) to prove that:
\begin{enumerate}
\item $H(X-Z)\le H(X-Y)+H(Y-Z)-H(Y)$
\item $H(X-Z)\le H(X+Y)+H(Y+Z)-H(Y)$
\end{enumerate}
[Hint: Use Part(b) with $X_1 = (X-Y,Y-Z)$, $X_2 = (X,Z)$, $X_{12} = (X,Y,Z)$ and $X_0 = X-Z$. ]

\item \underline{\it Sum and Difference of Entropy}  {\it (10 points)}\\
Use the previous part to conclude that for i.i.d $(Y,Y')$ random variables
\bea
H(Y-Y')-H(Y)\le 2(H(Y'+Y)-H(Y)).
\eea
\end{enumerate}

%%% Solution of  Problem 3: Entropy of Sum and Difference of I.I.D. Random Variables
%%%%%%%%%%%%%%%%%%%%%
\vskip 0.2in
\begin{solution}
{\bf Solution: }\\
\begin{enumerate}
\item 
\bea
I(X_1;X_2)&=&H(X_1)-H(X_1|X_2)\\
&\stackrel{(i)}{=}&H(X_1)-H(X_1|X_2,Y_2)\\
&\stackrel{(ii)}{\ge}&H(X_1)-H(X_1|Y_2)\\
&\stackrel{}{=}&I(Y_2;X_1)\\
&\stackrel{}{=}&H(Y_2)-H(Y_2|X_1)\\
&\stackrel{(iii)}{=}&H(Y_2)-H(Y_2|X_1,Y_1)\\
&\stackrel{(iv)}{\geq}&H(Y_2)-H(Y_2|Y_1)\\
&\stackrel{}{=}&I(Y_1;Y_2),
\eea
where
\begin{itemize}
\item[(i)] follows from the fact that $Y_2=G(X_2)$.
\item[(ii)] follows from conditioning reduces entropy.
\item[(iii)] follows from the fact that $Y_1=F(X_1)$.
\item[(iv)] follows from conditioning reduces entropy.
\end{itemize}
\item Clearly, $H(X_{12})\le H(X_1,X_2)$, thus $H(X_1)+H(X_2)-H(X_{12})\ge H(X_1)+H(X_2)-H(X_1,X_2)=I(X_1;X_2)$. Proof is completed by using (a) above.
\item
\begin{enumerate}
\item Let $X_1=(X-Y,Y-Z)$, $X_2=(X,Z)$, $X_0=X-Z$ and $X_{12}=(X,Y,Z)$. Thus we can have for some functions, F,G,R $X_0=F(X_1)=G(X_2)$ and $X_{12}=R(X_1,X_2)$. Using (b) above we have:
\bea
H(X,Y,Z)+H(X-Z)\le H(X-Y,Y-Z)+H(X,Z)
\eea
Rearranging and using independence and conditioning reduces entropy,
\bea
H(X-Z)\le H(X-Y) + H(Y-Z) - H(Y).
\eea
\item Replace $Y$ by $-Y$ and noting that $H(Y)=H(-Y)$ we have the result. 
\end{enumerate}
\item Use (c)-ii. for i.i.d. $X,Y,Z$ and using $H(Y)=H(X)$ and $H(X+Y)=H(Y+Z)$, we get,
\bea
H(X-Z)+H(X)\le 2H(X+Y)
\eea
Now replace $X=Y'$ and $Z=Y''$ to get the bound. 
\end{enumerate}
\end{solution}
\end{enumerate}


\begin{problem}
    \clearpage
    ~
    \clearpage
    ~
    \clearpage
    ~
    \clearpage
\end{problem}

\end{document}